How do language models represent a relational operation ---
\emph{currency-of}, \emph{capital-of} --- as distinct from the entity it
applies to? Across Qwen3-1.7B/8B and Gemma-2-9B, we find the operation
is carried by a \textbf{manipulable operator direction} in the residual
stream, largely \textbf{separable from the operand}: a two-way analysis
of variance factorizes the mid-network state as
\texttt{H\ \ensuremath{\approx}\ \ensuremath{\mu}\ +\ operand\ +\ operator\ +\ interaction},
\textasciitilde90\% additive, with the operator component dominant at
the query position. The direction is \textbf{causal}: adding
\texttt{v(op\_B)\ -\ v(op\_A)} flips the target-vs-source logit margin
for \textbf{all 20 ordered relation pairs in all three models} (a single
mid-workspace layer suffices), while matched-norm random controls do
nothing --- and a query-position activation patch reroutes the greedy
answer itself at the model\textquotesingle s competence ceiling. It
\textbf{generalizes}: directions built on half the operands flip the
held-out half; directions built in one wording of a relation transfer to
re-framed and fully re-lexicalized prompts; and the operation is
separable from its \textbf{surface realization} --- \emph{language-of}
and \emph{demonym-of} are distinct directions even though both emit
"Italian". The factorization is \textbf{specific to relational
retrieval}: under the identical pipeline, arithmetic and
comparison-logic operators show 2--4\ensuremath{\times} the interaction and fail held-out
generalization. Relational computation in these models is, to a first
approximation, an operand plus a transplantable operator.
